\section{文献综述}
\subsection{碳排放驱动因素研究}
现有研究关于碳排放驱动机制的讨论，主要沿着“恒等式分解—计量回归—系统协同”三条路径展开。第一类研究以Kaya恒等式和LMDI分解为代表，强调将碳排放变化拆解为人口、经济、能源效率和能源结构等可解释分量，优点是结构清晰、政策可读性强，适合做事后归因。第二类研究以STIRPAT扩展模型为代表，将IPAT关系随机化并引入弹性参数，使变量作用可在统计框架下估计，更适合考察多因素同时作用时的边际效应方向和强度。

近年来，驱动因素研究逐步从单一变量解释走向多机制协同识别。大量研究指出，煤炭依赖和高耗能产业比重通常与排放增长同向；而产业升级、能源替代、技术进步和治理投入在中长期具有减排潜力。值得关注的是，不同变量之间并非独立作用，例如能源结构变化往往伴随技术改造和产业调整，这使“单因素—单结论”的分析框架难以完整解释排放路径。因而，能够同时覆盖规模效应、结构效应和技术效应的综合识别框架，正在成为主流方向。

\subsection{碳排放预测方法研究}
碳排放预测方法已从传统统计模型扩展到机器学习和集成学习。传统时间序列模型（如AR类、趋势外推）在小样本、低频数据场景下通常更稳定，参数解释也较直接；但面对复杂非线性关系时，其刻画能力受到限制。机器学习模型在高维特征关系捕捉方面具备优势，尤其在样本容量较大时能体现更高拟合能力，不过其外推稳定性和可解释性在小样本条件下常受质疑。

为兼顾准确性与稳健性，越来越多研究采用多模型比较或模型集成策略。该思路不预设“复杂模型必然最优”，而是在统一评价指标下比较不同模型在外推任务中的表现，再据此选择适配模型。与此同时，SHAP等可解释方法被用于连接“黑箱预测”与“机制解释”，使预测模型在给出结果的同时能够提供变量重要性证据。总体来看，预测研究的发展趋势并非单纯追求复杂度，而是强调模型结构与数据规模、研究目标和解释需求之间的匹配。

\subsection{情景模拟与政策评估研究}
在达峰与减排研究中，情景模拟常用于评估政策强度变化对未来路径的影响。已有研究普遍构建基准情景与强化减排情景，通过设定能源结构、产业结构和技术进步等关键参数，比较不同路径下的峰值时点与峰值规模。该类方法的优势在于可将抽象政策目标转化为可量化路径，从而增强政策讨论的可操作性。

但现有情景研究也存在两个共性问题：一是部分参数设定与机制估计脱节，导致情景设计缺乏统计支撑；二是部分研究只给出终点数值，缺少对路径斜率、峰后下降速度和政策成本含义的解释。由此，情景模拟不仅需要“给出不同结果”，更需要解释“为什么出现差异、差异来自哪些变量变化”。这要求机制分析与情景分析在同一框架内衔接，而非各自独立。

\subsection{现有研究不足与本文定位}
综合已有文献，可以归纳出三点不足。其一，“机制解释”与“预测精度”常被分离处理，导致研究结论停留在解释层面或预测层面。其二，部分研究仍以样本内拟合表现评价预测能力，难以反映真实外推场景下的模型稳定性。其三，情景分析与机制识别衔接不足，政策建议容易停留于原则性表述，缺少与实证结果一一对应的证据链。

基于此，本文在国家级时间序列数据基础上，构建集扩展STIRPAT、EKC检验、中介路径分析、滚动一步预测比较与多情景模拟于一体的综合框架，力图实现三方面改进：第一，在同一数据口径下同时回答“为什么变化”和“如何变化”；第二，以滚动外推误差而非样本内拟合作为预测比较标准；第三，将机制证据嵌入情景路径解释，提升政策讨论的实证支撑力度。该定位使本文既保持统计解释的严谨性，也具备面向政策应用的可读性与可操作性。
